% Generated draft: requires mathematical and source review.
\documentclass[11pt]{article}
\usepackage[margin=25mm]{geometry}
\usepackage{fontspec}
\setmainfont{Latin Modern Roman}
\usepackage{amsmath,amssymb,mathtools,unicode-math}
\setmathfont{Latin Modern Math}
\usepackage{enumitem,needspace,xurl,hyperref}
\setlength{\parindent}{0pt}
\setlength{\parskip}{5pt}
\setlength{\emergencystretch}{3em}
\newcommand{\N}{\mathbb{N}}
\newcommand{\Z}{\mathbb{Z}}
\newcommand{\Q}{\mathbb{Q}}
\newcommand{\R}{\mathbb{R}}
\newcommand{\C}{\mathbb{C}}
\newcommand{\F}{\mathbb{F}}
\newcommand{\A}{\mathbb{A}}
\newcommand{\PP}{\mathbb{P}}
\newcommand{\E}{\mathbb{E}}
\newcommand{\eps}{\varepsilon}
\newcommand{\dd}{\,\mathrm{d}}
\begin{document}
\section*{UnsolvedMath ID: 30001619}
\textbf{Source problem code:} catalogue entry
\subsection{Definitions and mathematical statement}
A stochastic conservation law is a conservation equation incorporating randomness. A vanishing-viscosity perturbation introduces a viscosity parameter tending to zero. A dynamic large-deviation rate functional describes the exponential cost of deviations of trajectories in a specified asymptotic regime. Boundary terms here refer to contributions associated with the boundary in the corresponding large-deviation formulation; their precise meaning is not specified by the supplied record.

\textbf{Source question.} For the stochastic conservation law with vanishing-viscosity perturbation and dynamic large-deviation functional, can boundary conditions be added so that the corresponding boundary terms for dynamic large deviations are obtained?

\textbf{Unresolved formulation.} The supplied record does not specify the stochastic equation, spatial domain or dimension, time interval, unknown field, flux, viscosity operator, noise, initial conditions, admissible boundary conditions, or the dynamic large-deviation functional. It also does not specify the asymptotic scaling, trajectory space, topology, or the relation between the vanishing-viscosity limit and the large-deviation limit. Consequently, the question asks for a boundary extension of an unidentified underlying formulation. No particular equation, class of boundary data, formula for a boundary contribution, or universal quantifier over conservation laws can be supplied without additional source information. A precise theorem-level target therefore remains undetermined.
\subsection{Short English statement}
The question asks whether boundary conditions can be incorporated into the intended stochastic conservation law with vanishing viscosity, and whether the associated dynamic large-deviation boundary contributions can be derived. The supplied record does not identify the equation or large-deviation framework, so it supports this research question but does not determine a precise theorem statement.
\subsection{Sources}
\begin{itemize}
\item Supplied local UnsolvedMath registry record, permanent ID \(30001619\), title \emph{Boundary Terms in Dynamical Large Deviations}, exact source problem code \texttt{catalogue entry}, in the collection \emph{Oberwolfach Reports Open Problems}. No individual proposer is supplied; the proposed year is recorded as 2010.
\item Supplied report citation: \emph{Large Scale Stochastic Dynamics} (2011), Oberwolfach Reports, pp. 2877--2975, DOI 10.4171/owr/2010/50. Supplied literature-source label: ``original OWR report.'' \url{https://doi.org/10.4171/owr/2010/50}. The record describes Oberwolfach Reports as documenting mathematical workshops and their open-problem discussions. No problem-specific page or passage is supplied.
\item Supplied literature-source label: ``Bertini--Landim--Mourragui.'' \url{https://arxiv.org/abs/0804.2458}. No paper title, theorem reference, or precise connection to the intended stochastic equation is supplied.
\item Supplied literature-source label: ``Franco--Goncalves--Landim--Neumann.'' \url{https://arxiv.org/abs/2203.14417}. No paper title, theorem reference, or precise connection to the intended stochastic equation is supplied.
\item Literature assessment dated exactly \texttt{2026-08-21}: the supplied curation labels the problem ``partially solved'' and states that dynamic large-deviation principles with boundary conditions have been proved for several one-dimensional boundary-driven exclusion models, including weakly asymmetric and Robin-boundary cases, while the general stochastic-equation and boundary-term question remains open. This is supplied dated curation derived from the OWR record, not an independently verified mathematical or literature assessment; its time-sensitive status and cited claims require independent verification.
\end{itemize}
\end{document}
